\chapter{问题定义}
\label{chap:setup}

本章形式化本文研究的问题。我们首先给出基本设定与 GRPO 的优势估计，然后指出在多轮智能体场景下该估计的两类信号污染，最后把问题重新表述为两阶段形式并给出持续学习的目标度量。

\section{基本设定}
\label{sec:setup-basic}

\textbf{策略与算法。} 基座模型为 Qwen3.6-27B，采用 GRPO（无 critic）作为强化学习算法。对每个查询，策略 rollout 出 $M=8$ 条轨迹，组内归一化得到每条样本的优势：
\begin{equation}
\label{eq:grpo}
A_i = \frac{r_i - \mathrm{mean}(r)}{\mathrm{std}(r) + \epsilon},
\end{equation}
其中 $r_i$ 为第 $i$ 条轨迹的奖励，$\mathrm{mean}(r)$、$\mathrm{std}(r)$ 为组内均值与标准差，$\epsilon$ 为数值稳定小常数。

\textbf{数据与任务。} 评测与训练覆盖纯文本 195 个任务，跨 9 个能力域（详见第~\ref{chap:exp} 章）。任务以多轮会话形式组织：每个会话含多条顺序查询，共享上下文。每个任务含一份声明（首条真实查询 $q_1$、隐藏目标、判分准则、用户画像）与一份初始文件系统（沙箱种子）。

\textbf{目标。} 持续学习的目标是最大化综合指标
\begin{equation}
\label{eq:clscore}
\mathrm{CL\ Score} = \mathrm{New\ Task\ Perf} - \alpha \cdot \mathrm{Old\ Task\ Forgetting},
\end{equation}
其中 $\alpha=1.0$，New Task Perf 衡量新任务学习效果，Old Task Forgetting 衡量旧任务评测分数相对上一阶段的下降量。

\section{GRPO 优势估计的隐含前提}
\label{sec:setup-premise}

贯穿全文的一个隐含前提是：式~\eqref{eq:grpo} 的组内归一化优势估计\textbf{仅当} $M$ 条轨迹的差异纯粹来自策略采样的随机性时才无偏。换言之，组内方差应当只反映"同一策略在同一查询下的采样随机性"，而不应混入环境噪声或数据前提错位。任何对这一前提的违反都会把与策略无关的方差注入 $A_i$，从源头污染学习信号。第~\ref{chap:method} 章的 \S\ref{sec:rollout}--\ref{sec:usersim} 即从源头守住这一前提。

\section{多轮智能体场景的两类信号污染}
\label{sec:setup-corruption}

在多轮智能体强化学习中，上述前提并不自动成立。本文识别出两类在源头破坏优势估计的污染。

\subsection{组内环境发散}
\label{sec:setup-divergence}

跨多查询会话，同一组的 $M$ 个执行环境会逐渐分叉。设第 $k$ 条查询处理后，第 $j$ 个槽位的环境状态为 $e_k^{(j)}$。若不在查询边界对齐各槽位，则第 $k{+}1$ 条查询的 $M$ 个槽位将从不同状态 $\{e_k^{(1)},\dots,e_k^{(M)}\}$ 出发，组内比较就混入了历史路径差异：某个槽位可能并非因策略更差、而是因第 $k$ 条查询留下的不利状态被惩罚。于是优势 $A_i$ 吸收了随查询序号 $k$ 累积的环境方差，且组内标准差 $\mathrm{std}(r)$ 会被这一环境噪声放大，导致信用分配（credit assignment）方向出错。

\subsection{多轮前提漂移}
\label{sec:setup-premise-drift}

后续查询 $q_{k+1}$ 通常依赖上一轮执行后的状态 $e_k$（例如"这份 PPT 第 3 页数据错了"）。记 $\phi(q_{k+1})$ 为 $q_{k+1}$ 的前提函数，当状态 $e_k$ 满足该前提时 $\phi(q_{k+1})(e_k)$ 为真。若 $q_{k+1}$ 在数据采集时被静态写死，而训练时 $e_k$ 是策略 $\pi_\theta$ 的随机输出，则
\begin{equation}
\label{eq:premise}
\Pr\big[\phi(q_{k+1})(e_k)\big] < 1,
\end{equation}
且该概率随策略变强而\textbf{持续漂移}。前提失败会产生污染信号：模型要么因对不存在的问题硬编一个修复而被奖励，要么因如实指出该问题不存在而被惩罚；同时，前提不成立的样本会把环境噪声注入组内归一化，稀释信号。

\textbf{关键澄清。} 即便回流数据里含有真实的后续查询 $q_2,\dots,q_K$，它们同样是在采集时针对\textit{原始}会话的执行结果写下的；训练 rollout 中策略对 $q_1$ 产出的状态与该原始结果不同，这些真实后续查询的前提同样不成立。直接复用它们会重新引入前提漂移——这正是本文只保留种子 $q_1$、在线构造后续轮次的原因。

\section{两阶段重新表述}
\label{sec:setup-reframe}

被污染的优势在下游无法修复：任何回放加权或 KL 锚定都救不回一个符号本就错了的梯度。这把持续智能体强化学习重新表述为一个\textbf{两阶段}问题：

\begin{enumerate}
  \item \textbf{产出干净信号}（第~\ref{chap:method} 章 \S\ref{sec:rollout}--\ref{sec:usersim}）：用沙箱 winner 同步会话调度消除组内环境发散（\S\ref{sec:setup-divergence}），用三智能体在线构造多轮查询消除前提漂移（\S\ref{sec:setup-premise-drift}），使式~\eqref{eq:grpo} 的优势估计恢复无偏。
  \item \textbf{巩固干净信号}（第~\ref{chap:method} 章 \S\ref{sec:obj}--\ref{sec:reweight}）：在干净信号之上施加持续学习目标、抗遗忘回放缓冲区与块级重加权，在不遗忘的前提下巩固学习成果。
\end{enumerate}

本文的整个系统围绕这一两阶段主旨组织：没有任何一个组件被当作独立技巧来呈现，每一项都是服务于"先产出、再巩固一个未被污染信号"这一主旨的手段。
